\section{完整支持的用途：来源日志诊断}\label{sec:task}
客户端接收一个固定事件、固定读值的输入，以及不完整来源／写序日志 $\alpha$。一个 RF 字段记录某个读的来源身份，一个 CO 字段记录某写对的方向；未记录的字段保持自由。这里“不完整”不表示可以缺少定理所需的阶段事件或等待读值。

对一致日志，客户端返回每个读全部仍可行的来源，以及每个 CO 对仍可行的方向。它们是逐项可能性，不能任意组合为同一个执行。对冲突日志，客户端逐项尝试删除记录，仅在删除后仍 UNSAT 时保留删除，得到\emph{删除意义下极小}的冲突集。这不等于最小基数冲突，也不要求枚举全部冲突。

单字段修复使用完整支持。对每个记录字段 $j$，暂时撤销它；若其余日志仍冲突，修改此字段无法修复。否则查询剩余日志的完整支持：RF 字段返回该读所有有支持的替代来源，CO 字段返回有支持的相反方向。移除原值后，即得到恰好修改一个字段的所有修复。

\begin{cprop}[单字段修复的完备性]\label{prop:repair}
上述过程返回且仅返回能够使整个日志一致的单字段替换。
\end{cprop}
\begin{cproof}
任何修改字段 $j$ 的合法修复都满足撤销 $j$ 后的其余日志，所以新值必在该次支持中。反之，每个返回支持都有满足其余日志和该新值的完整合法见证。遍历全部字段故既不漏修复，也不引入虚假修复。
\end{cproof}

第~\ref{sec:motivation}节的冲突 $\{d,\ell,\neg q\}$ 是删除极小的，并有恰好三种单字段修复：将数据来源从 B 第二轮自身的 $y$ 写改成 A 第二轮的 $y$ 写；将同步来源从 A 第一轮晚发布改成早发布；或将写序方向从 $\neg q$ 改为 $q$。三个选项各有合法见证，它们作为替代建议分别成立，不承诺任意联合使用仍成立。

一次 SAT 调用只能给出其中一个模型，不能用该模型没有选中的来源证明其他来源不可能。SAT 也能实现相同任务，但需要按其查询接口覆盖所有支持。本文同时使用逐原始位完整支持的 SAT 客户端，以及按当前操作只求一致性或相关字段取值的按需 SAT 客户端；两者均返回相同的诊断结果。支持编译的价值在于一次构造后重复使用精确表示，而非将“完整支持”误写成 SAT 无法完成的功能。

实验日志由交换内核生成，覆盖早／晚来源、CO 倒置、冲突及撤销。$n$ 轮任务产生 $8n-7$ 份诊断报告、$20n-19$ 次支持接口调用和 $12n-11$ 个不同假设集合；不使用查询结果缓存。客户端本身不依赖族生成器，可接受任意符合语义输入格式的固定片段。本文尚未用工业部署采集的日志评测，因此将用途定位为可执行的诊断任务，而非部署效果证明。
